% Introduction

Reinforcement learning for control aims to learn decision-making policies through interaction with dynamic environments. In such settings, it is not sufficient for a method to perform well in a single run; evaluation should also be reproducible, comparable, and robust across environments, training budgets, and random seeds. This project investigates how Group Relative Policy Optimization (GRPO), originally developed in the context of language-model reinforcement learning and mathematical reasoning, can be adapted and evaluated in control-oriented benchmark environments \citep{shao2024deepseekmath}.

The project uses PPO, SAC, and TD3 as baseline algorithms. These methods represent central approaches in modern deep reinforcement learning: PPO is an on-policy policy-gradient method based on stabilized policy updates, SAC is an off-policy maximum-entropy actor-critic method, and TD3 is a deterministic actor-critic method designed for continuous action spaces \citep{schulman2017proximal,haarnoja2018sacapps,fujimoto2018td3}. In this report, these algorithms are not the proposed methodological contribution. Instead, they define the comparison surface needed to evaluate a later GRPO-based control method.

GRPO is relevant because it constructs policy updates from group-relative signals rather than relying solely on a conventional value-function baseline. However, transferring this idea to control is not direct. Control environments require sequential credit assignment, handling of continuous or structured action spaces, and stable learning from trajectories whose outcomes may depend strongly on the environment, seed, and training budget. A GRPO variant for control therefore requires both a precise formal problem formulation and a reliable baseline foundation.

At the midway stage, this report does not present a completed GRPO-control algorithm. The contribution is instead to establish the experimental and methodological foundation for the final comparison. The report introduces the Markov decision process notation, documents the baseline implementation, and validates a complete experimental matrix with three algorithms, three environments, and five seeds, corresponding to 45 runs and 900 evaluation rows. The vector-based control environments are run through ObjectRL, while CarRacing is handled by project-side CNN implementations because the image observations used in this setup do not fit directly into ObjectRL's standard vector-observation workflow \citep{baykal2025objectrl,towers2024gymnasium}.

The results should therefore be interpreted as a midway foundation rather than as final performance conclusions. No hyperparameter optimization was performed, and the training budgets were chosen to validate the pipeline and comparison setup rather than to establish asymptotic performance. The results primarily show that the environment adapters, training scripts, evaluation protocol, result files, and figure generation work end-to-end. This provides the necessary basis for the later GRPO evaluation.

The remainder of the report is structured as follows. Section~\ref{sec:related} reviews related work on PPO, SAC, TD3, and GRPO. Section~\ref{sec:methodology} describes the formal MDP notation and experimental methodology. Section~\ref{sec:theory} outlines the theoretical and notational basis for the later GRPO extension. Section~\ref{sec:experiments} presents the completed midway baseline experiments. Finally, Section~\ref{sec:conclusion} summarizes the current limitations and the next steps toward the final GRPO-based control method.